\section{\texorpdfstring{$G_{5,14}$}{G5,14} の五node単位枝の無上界排除}
\label{ug25:section}

本節では二つの単位枝を同時に用いる．単位枝一つの三角形条件だけでは残った
$\{5,0,0\}$ 配分が，もう一つの単位枝と正の三高 core の行和を合わせると
排除されることを示す．分母や行列成分に探索上限を置かない．

\subsection{仮定と，利用する局所事実}

五つの cores を $0,1,2,3,4$ とし，
$q_{01}=q_{23}=1$，他の八分母は一より大きい奇数とする．
core $i$ の ray $j$ から $k$ への整数成分を $p^i_{jk}$ と書き，
元の条件を
\begin{align}
 \sum_{k\ne i,j}p^i_{jk}&=-q_{ij},&
 p^i_{jk}p^i_{kj}&=q_{ij}q_{ik}-1,\label{ug25:local}\\
 w^i_k-\frac{p^i_{jk}}{q_{ij}}w^i_j&>0\quad(w^i_j>0),&
 n_{ij;k}&=1-\frac{p^i_{jk}+p^j_{ik}}{q_{ij}}\in\mathbb Z_{\ge0},
 \label{ug25:positive-node}\\
 \sum_{k\ne i,j}n_{ij;k}&=5,&
 D^i_j+D^j_i&=0,\qquad
 D^i_j=-q_{ij}-\sum_{k\ne i,j}p^i_{kj}
 \label{ug25:budget-D}
\end{align}
とする．本節の結論はこの標準五次接着の必要条件に対するものであり，
任意の Lagrangian torus fibration の分類を主張しない．

正の三高 core の正 pair graph は完全マッチングである
（三高 matching 定理~\ref{mh18:matching}）．従って唯一の低行は
正一成分・負二成分を持ち，その二つの負位置同士も正 pair を作る．
以下ではこの帰結を明示的に使う．四高 core の正 graph の分類は使わない．

単位枝の両低行の正位置は異なる．同じ位置ならその node が負になるからである．
正位置を $a,b$，共通の負位置を $c$ として
\[
 \begin{array}{c|ccc}
 &a&b&c\\\hline
 i&u&-t&-s\\
 j&-j_0&v&-s'
 \end{array},\qquad
 u=s+t-1,\quad v=s'+j_0-1
\]
と書くと，$n_{ij;c}=1+s+s'$，従って $s+s'\le4$ である．
共通負位置の低 column を $-S,-T$ と書けば
$q_{ic}=sS+1$，$q_{jc}=s'T+1$ である．
高枝の node から得る $\alpha,\beta\in\{-4,-3,-2,-1,0,1\}$ により，
外側 core $c$ の pair equation は
\begin{equation}\label{ug25:negative-pair}
 ((s\alpha+1)S+\alpha)((s'\beta+1)T+\beta)
 =(sS+1)(s'T+1)-1
\end{equation}
となる．ここから得る必要条件を再掲する：
\begin{center}\small
\begin{tabular}{c|c|c|c}
外側 pair&$s+s'$&$(q_{ic},q_{jc})$&必要な低成分\\\hline
正&3&$(2R+1,3)$&$(s,s',S,T)=(2,1,R,2)$\\
正&3&$(3,2R+1)$&$(s,s',S,T)=(1,2,2,R)$\\
負&3&$(3,3)$&---\\
負&4&$(3,7)$&$(1,3,2,2)$ または $(2,2,1,3)$\\
負&4&$(5,5)$&$(2,2,2,2)$\\
負&4&$(7,3)$&$(3,1,2,2)$ または $(2,2,3,1)$
\end{tabular}
\end{center}
$R$ は任意の正整数である．特に全配分は $\{4,1,0\}$ または
$\{5,0,0\}$ であり，五nodeの場合には $S+T=4$ である．
この表は補題~\ref{ut24:negative-classification} の全解から従うが，
必要な無上界性は次の因数分解でも確認できる．
負 pair に対して $\alpha=-A,\beta=-B$ とすると，
\eqref{ug25:negative-pair} の左辺引く右辺は
\[
 CST+DS+ET+AB,
 \quad C=ss'(AB-1)-sA-s'B+1,
 \quad D=s(AB-1)-B,
 \quad E=s'(AB-1)-A.
\]
$A,B\ge2$ は係数の非負性で，$A=B=1$ は狭義負で除かれる．
交換により $A=1,2\le B\le4$ としてよい．
$s=1$ なら $S=B-1-(T-1)/(s'T+1)$ より $T=1$，
$s=2,B=2$ なら $T(S-s'+1)=2$，
$s=3,s'=1,B=2$ なら $S(T-1)=2$ である．
$s\ge2,B\ge3$ は係数の符号で除かれる．奇数性が上表の全点を与える．
正 pair では $\alpha,\beta\in\{0,1\}$ となり，
$(1,0)$ が $s'=1,T=s$，$(0,1)$ がその交換を与える．
奇数性と $s+s'\le4$ が二つの正族を与える．

次に反対符号位置の低行を $v,-j$，低 column を $V,-J$ とする．
$n$ を単位枝上のその位置の node とすると $j=v+n-1$ である．
外側 pair equation は
\begin{equation}\label{ug25:opposite-pair}
 ((v\alpha-1)V+\alpha)((j\beta+1)J+\beta)
 =(vV+1)(jJ+1)-1.
\end{equation}
本節で必要な負 pair，$n\in\{0,1\}$ の全族は
\begin{center}\small
\begin{tabular}{c|c|c|c|c|c}
$n$&$v$&$j$&$V$&$J$&$(\alpha,\beta)$\\\hline
0&2&1&$2r$&$2r$&$(0,-3)$\\
1&正の偶数&$v$&任意の正整数&1&$(-1,-1)$\\
1&1&1&$2r$&$2r$&$(0,-2)$\\
1&2&2&1&1&$(0,-3)$\\
1&1&1&$2r$&$k\in\{2,4\}$&$(-k,-1)$
\end{tabular}
\end{center}
であり，$r\ge1$ である．正 pair，$n=0$ の全族は
\begin{equation}\label{ug25:positive-zero}
 \alpha=\beta=1,\qquad j=v-1,\qquad J=V-1,
 \qquad v,V\ge2,\quad v\not\equiv V\pmod2
\end{equation}
である．これらも命題~\ref{ut24:opposite-classification} の一部だが，
本節で使う範囲の完全性を短く確認しておく．負 pair では
$\alpha=-A,\beta=-B$，$0\le A\le4,1\le B\le4$ と置く．差は
\[
 [vj(AB-1)+jB-vA-1]VJ
 +[v(AB-1)+B]V+[j(AB-1)-A]J+AB.
\]
$A\ge1,B\ge2$ は係数の非負性で除かれる．
$A=B=1$ は $(n-2)VJ+V-J+1=0$ なので，$n=1,J=1$ となる．
$A>1,B=1$ では $J$ 係数が負でなければならず $j=1$ であり，
$n=0,v=2$ は整数解なし，$n=1,v=1$ は $J=A$ である．
$A=0$ では $d=B-v$ と置くと
\[
 J((jd-1)V-j)=-dV.
\]
$d\le0$ は不可能，$jd=1$ は $j=d=1,J=V$，
$jd>1$ は $d=V=J=1$ となる．これに奇数性を課すと上表を得る．
正 pair では $(\alpha,\beta)=(1,1)$ 以外が排除され，
$J(1-nV)=V-1$ から \eqref{ug25:positive-zero} を得る．

\subsection{五node仮定の正規形}

\begin{lemma}\label{ug25:normal-form}
単位枝 $01$ が五nodeを持つなら，標識を取り直してその位置を $2$ とし，
低 row と低 column を，列順 $(3,4,2)$ で
\begin{equation}\label{ug25:normal-rows}
 \begin{array}{c|c|c}
 &\text{低 row}&\text{低 column}\\\hline
 0&(u,-t,-s)&(A,1-B,-S)\\
 1&(1-u,t+1,-s')&(1-A,B,-T)
 \end{array}
\end{equation}
と書ける．ここで
\[
 u=s+t-1,\quad s+s'=S+T=4,\quad u,A,B\ge2,
 \quad u\not\equiv A\pmod2,
\]
かつ $(s,s',S,T)$ は上の五node表に属する．従って
\begin{align}
 q_{02}&=sS+1\in\{3,5,7\},&q_{12}&=s'T+1,\label{ug25:qnormal1}\\
 q_{03}&=uA+1\ge7,&q_{13}&=(u-1)(A-1)+1,\label{ug25:qnormal2}\\
 q_{04}&=t(B-1)+1,&q_{14}&=(t+1)B+1.\label{ug25:qnormal3}
\end{align}
core $0$ の pair $23$ は負，core $1$ の pair $23$ は正である．
core $3$ の低行の正位置は $4$ である．
\end{lemma}

\begin{proof}
五nodeの他の二位置は node 零である．
\eqref{ug25:opposite-pair} の $n=0$ の二族により，
正側低 column と負側低 column の絶対値の差は零または一である．
五node位置では $S+T=4$ であり，対向低 $D$ の相殺は全低 column の和を
$-2$ とするので，残る二差の和は二である．従って両方の差は一で，
両外側 pair は正となる．

五node位置が $4$ なら，core $2,3$ の pair $01$ がともに正となる．
matching によりそれぞれの低行の正位置が $4$ となり，単位枝 $23$ の
位置 $4$ の node が負になる．従って五node位置は $2$ または $3$ である．
交換して $2$ とし，$0,1$ を交換して core $0$ の低正位置を $3$，
core $1$ の低正位置を $4$ とする．
\eqref{ug25:positive-zero} を二つの位置に適用すると
\eqref{ug25:normal-rows} と分母式を得る．
$u,A\ge2$ は異なる偶奇なので $uA\ge6$ である．
matching が残りの符号を決め，core $3$ で pair $01$ が正であることから
その低正位置は $4$ となる．
\end{proof}

\subsection{もう一方の単位枝による二つの還元}

core $2$ で pair $01$ は負であるため，matching によりその低正位置は
$0$ または $1$ である．$4$ なら pair $01$ が正になるので不可能である．

\begin{lemma}\label{ug25:case-A}
core $2$ の低正位置が $0$ なら
\[
 (q_{02},q_{03},q_{12},q_{13})=(7,7,3,3),
\]
かつ $(s,s',S,T,u,A,t)$ は
\begin{equation}\label{ug25:A-cases}
 (3,1,2,2,3,2,1),\quad
 (2,2,3,1,2,3,1),\quad
 (2,2,3,1,3,2,2)
\end{equation}
の三つのいずれかである．
\end{lemma}

\begin{proof}
単位枝 $23$ の両低正位置は $0,4$ なので，共通負位置は $1$ である．
core $1$ の pair $23$ が正であるから，その位置の node は四であり，
位置 $0$ の node $n$ は零または一である．
core $0$ の pair $23$ は負である．その反対符号位置の表を用い，
core $2$ 側を正の低行 $v$，低 column $V$，core $3$ 側を
負の低行 $-j$，低 column $-J$ とする．
\[
 q_{02}=vV+1\le7,\qquad q_{03}=jJ+1\ge7.
\]
$n=0$ の例外は $v=2,j=1,V=J=2r$ であり，第一不等式から $r=1$，
従って $q_{03}=3$ となり矛盾する．
$n=1$ の通常負族は $j=v,J=1$ であり，
$v+1\le vV+1\le7\le v+1$ から $v=6,V=1$ となる．
例外 $v=j=1,V=J=2r$ も両分母を七に強制する．
他の二例外では $q_{03}\le5$ なので不可能である．
従って常に $q_{02}=q_{03}=7$ である．
最初の単位枝の五node表から $q_{12}=3$，
$uA=6$ から $(u,A)=(2,3),(3,2)$ かつ $q_{13}=3$ を得る．
五node表と $t=u-s+1\ge1$ を組み合わせると
\eqref{ug25:A-cases} が尽くされる．
\end{proof}

\begin{lemma}\label{ug25:case-B}
core $2$ の低正位置が $1$ なら
\[
 (q_{02},q_{03},q_{12},q_{13})=(3,7,7,3),
\]
かつ $(s,s',S,T,u,A,t)$ は
\begin{equation}\label{ug25:B-cases}
 \begin{split}
 &(1,3,2,2,2,3,2),\quad (1,3,2,2,3,2,3),\\
 &(2,2,1,3,2,3,1),\quad (2,2,1,3,3,2,2)
 \end{split}
\end{equation}
の四つのいずれかである．
\end{lemma}

\begin{proof}
単位枝 $23$ の共通負位置は $0$ で，そこでの pair $23$ は負である．
共通負位置の表により $(q_{02},q_{03})$ は
$(3,3),(3,7),(5,5),(7,3)$ のいずれかである．
$q_{03}\ge7$ により $(3,7)$ だけが残る．
最初の単位枝の五node表から $q_{12}=7$，また $uA=6$ から
$q_{13}=3$ を得る．同じ表の $(3,7)$ の二つの低成分と
$(u,A)=(2,3),(3,2)$ を組み合わせると
\eqref{ug25:B-cases} を得る．
\end{proof}

\subsection{二つの固定小分母を持つ core の有限因子障害}

ここまで残った core $1$ の低 row と低 column を，列順 $(3,2,4)$ で
\begin{equation}\label{ug25:fixed-core}
 (-(u-1),-s',v),\qquad (-(A-1),-T,B),\qquad v=t+1
\end{equation}
と書く．最初の二分母を $m,n$ とすると，七つのケースは次表になる．
$a=m-(A-1)$，$b=n-T$ とした．
\begin{center}\small
\begin{tabular}{c|cc|ccc|c}
ケース&$m$&$n$&$v$&$a$&$b$&$B$ の当初の条件\\\hline
$A_1$&3&3&2&2&1&$B\ge2$\\
$A_2$&3&3&2&1&2&$B\ge2$\\
$A_3$&3&3&3&2&2&$B\ge2$\\\hline
$B_1$&3&7&3&1&5&$B\ge2$\\
$B_2$&3&7&4&2&5&$B\ge2$\\
$B_3$&3&7&2&1&4&$B\ge2$\\
$B_4$&3&7&3&2&4&$B\ge2$
\end{tabular}
\end{center}

\begin{lemma}[固定二分母の因子判定]\label{ug25:factor-test}
\eqref{ug25:fixed-core} を持つ正の三高整数 core が存在するなら，
ある正整数 $L,M$ に対して
\begin{equation}\label{ug25:factor-equations}
 LM=mn-1,\qquad
 \frac{mvB+m-1}{L+a},\quad
 \frac{nvB+n-1}{M+b}\quad\text{はいずれも正整数},
\end{equation}
かつ
\begin{equation}\label{ug25:linear-equation}
 CB+E=0,
 \quad
 \begin{cases}
 C=mv(M+b)+nv(L+a)-(v+1)(L+a)(M+b),\\
 E=(m-1)(M+b)+(n-1)(L+a)-(L+a)(M+b)
 \end{cases}
\end{equation}
が成り立つ．これは任意に大きい $B$ に対する必要条件である．
\end{lemma}

\begin{proof}
matching により，二つの負低行位置の間の pair は正である．
その二 entries を $L,M>0$ とすれば $LM=mn-1$ である．
当該二行の残りの高 entry は行和から $-(L+a),-(M+b)$ である．
分母 $vB+1$ の ray からの逆 entries は pair equations により
\eqref{ug25:factor-equations} の二商の負数となる．
それらの絶対値の和は最後の行和から $(v+1)B+1$ である．
分母を払うと \eqref{ug25:linear-equation} を得る．
整除条件を落としていない点が重要である．
\end{proof}

$A$ の三ケースでは $LM=8$，$B$ の四ケースでは $LM=20$ である．
全正因子を代入して得る $B=-E/C$ は次の通りである．
\begin{center}\small
\begin{tabular}{c|rrrr}
&$L=1$&$L=2$&$L=4$&$L=8$\\\hline
$A_1$&$-1/3$&$-1/3$&$\ast$&$-1/3$\\
$A_2$&$-1/3$&$\ast$&$-1/3$&$-1/3$\\
$A_3$&$-4/3$&$-2/3$&$-2/3$&$-4/3$
\end{tabular}
\end{center}
$\ast$ では $C=E=0$ だが，二商の一方が
$(6B+2)/6=B+1/3$ となり，整数性に反する．
\begin{center}\small
\begin{tabular}{c|rrrrrr}
&$L=1$&$L=2$&$L=4$&$L=5$&$L=10$&$L=20$\\\hline
$B_1$&$-12/67$&$-1/6$&0&0&$3/14$&$4/3$\\
$B_2$&$7/9$&$-3/4$&$-1/3$&$-3/11$&$\dagger$&$-3/7$\\
$B_3$&$-3/7$&$\dagger$&$3/11$&$1/3$&$3/2$&$-31/9$\\
$B_4$&$-2/3$&$-2/7$&0&$2/5$&$-2/3$&$-32/67$
\end{tabular}
\end{center}
$\dagger$ では $C=0,E\ne0$ である．従ってどの列にも正整数 $B$ はない．
この三十六場合は固定された八と二十の全正因子であり，有限範囲の数値探索ではない．

\begin{theorem}[$G_{5,14}$ の五node排除]\label{ug25:no-five}
\eqref{ug25:local}--\eqref{ug25:budget-D} を満たす標準五次接着で，
高分母 support が $G_{5,14}=K_5\setminus\{01,23\}$ ならば，
二つの単位枝の node 配分はともに $\{4,1,0\}$ である．
特にどの単位枝も $\{5,0,0\}$ を持たない．
\end{theorem}

\begin{proof}
単位枝 $01$ が五nodeを持つと仮定する．
補題~\ref{ug25:normal-form} で正規化し，core $2$ の低正位置により
補題~\ref{ug25:case-A} または補題~\ref{ug25:case-B} を適用する．
残る七つの局所型は全て補題~\ref{ug25:factor-test} と全因子表に反する．
単位枝 $23$ については二つの枝を交換すれば同じである．
一般の単位枝配分は既に二種類に限定されているので結論を得る．
\end{proof}

\begin{corollary}[残る配置の小分母]\label{ug25:small-three}
$G_{5,14}$ の各単位枝について，共通負位置の二つの隣接高分母の少なくとも
一つは三である．その外側 pair が負なら両方とも三であり，正なら
$(3,2R+1)$ または $(2R+1,3)$，$R\ge1$ である．
\end{corollary}
\begin{proof}
五nodeが排除された後の共通負位置の表を読むだけでよい．
\end{proof}

\paragraph{検証と残る問題．}
\nolinkurl{verify_g514_five_node_v25.py} は旧 verifier を import せず，
必要な共通負位置の有限係数分類，七局所型の再構成，三十六正因子の
整数係数線形式と整除障害を Python 標準ライブラリの整数・有理数で確認する．
無限尺度を除くのは式と整除障害であり，有界 scan の不在証明への外挿ではない．
本定理は $G_{5,14}$ support 自体を排除していない．
二つの $\{4,1,0\}$ 配分の同時接着，$G_{5,16}$ の五node sector，
四高 core の完全 matching 型における整数模型の分類と大域整合性は未解決である．
